% Volume VIII: Applications and Case Structures
% Part XV. Domains of Difficult Judgment
% Chapter 86: Medicine and Scarce Data
% Spine: proof-spines/ch086-spine.md

\chapter{Medicine and Scarce Data}
\label{ch:ch086}

In medicine, good judgment often depends on learning carefully from data too scarce, noisy, and heterogeneous to bear the weight we want to place on it. After the previous chapter's treatment of high-stakes public evidence, this one turns to the clinic and laboratory, where rare conditions, small cohorts, noisy imaging, and shifting baselines make overconfident updating especially attractive and especially costly. It will use medical-imaging and model-training examples to show why constrained forms of revision can outperform wholesale retraining when evidence is limited: the problem is not merely lack of data, but lack of permission to generalize freely from it. The chapter then connects bedside decision, trial design, and algorithmic assistance under a single norm of epistemic thrift, asking what clinicians may infer, what they must bracket, and how uncertainty should be communicated to patients. Its larger point is that skepticism here is not therapeutic defeatism but a method for preserving signal when the environment punishes exuberant inference.
